\uuid{f15u}
\exo7id{1656}
\titre{exo7 1656}
\auteur{legall}
\organisation{exo7}
\datecreate{1998-09-01}
\isIndication{false}
\isCorrection{false}
\chapitre{Réduction d'endomorphisme, polynôme annulateur}
\sousChapitre{Diagonalisation}
\cpge{mpsi-mp2i.dimension-finie.a, pcsi.dimension-finie.a, ptsi.dimension-finie.a, mpsi-mp2i.dimension-finie.b, pcsi.dimension-finie.b, ptsi.dimension-finie.b, bcpst1.espace-vectoriel-kn.b, mp-mpi.reduction-endomorphismes.d, pc.reduction.c, psi.reduction-endomorphismes.c, pt.reduction-endomorphismes.c, mp-mpi.reduction-endomorphismes.g, pc.complements-espaces-vectoriels.d, psi.complements-espaces-vectoriels.d, mp-mpi.reduction-endomorphismes.i, pc.reduction.d, psi.reduction-endomorphismes.d, bcpst2.valeurs-propres-vecteurs-propres.b}
\module{Algèbre}
\niveau{L2}
\difficulte{}

\contenu{
\question{
Soit $  E  $ un $  \mathbb{K} $-espace vectoriel de dimension finie et $  f
\in \mathcal{L} (E)  $ diagonalisable. Montrer que les propri\'et\'es suivantes sont \'equivalentes~:
\vskip1mm
(i) La famille $  \{ id, f, f^2, \ldots , f^{n-1}\}   $ est libre.
\vskip1mm
(ii) Il existe $  x \in E   :   \{ x, f(x), f^2(x), \ldots , f^{n-1}(x)\}   $ engendre $  E  .$
\vskip1mm
(iii) Les valeurs propres de $  f  $ sont simples.
}
}
